\section{Introduction}

\begin{frame}{Problem Formulation: Face Reenactment}
    \begin{columns}[c]
        \begin{column}{0.54\textwidth}
            \small
            \textbf{Core Objective}
            \begin{itemize}
                \item Every method works with two roles.
                \item The \textbf{identity image} supplies the face.
                \item The \textbf{motion source} supplies the movement.
                \item Each method uses the pair in its own way.
            \end{itemize}

            \vspace{0.15cm}
            \textbf{Key Technical Challenges}
            \begin{itemize}
                \item \textbf{Identity preservation:} Avoid warping face features into driving face.
                \item \textbf{Non-rigid deformation:} Facial skin stretches (not just rigid 2D shift).
                \item \textbf{Occlusion handling:} Mouth opening reveals unseen areas (teeth, tongue).
            \end{itemize}
        \end{column}

        \begin{column}{0.44\textwidth}
            \begin{block}{Input / Output Setup}
                \small
                \textbf{Inputs:}
                \begin{itemize}
                    \item Identity image: supplies the face.
                    \item Motion source: supplies the movement.
                \end{itemize}
                \vspace{0.1cm}
                \textit{Each method gives the inputs their own names.}
                \vspace{0.1cm}
                \textbf{Output:}
                \begin{itemize}
                    \item Reenacted video of the identity face
                    \item Mimicking expressions and head pose of the motion source
                \end{itemize}
            \end{block}
        \end{column}
    \end{columns}
\end{frame}
